% Generated draft: requires mathematical and source review.
\documentclass[11pt]{article}
\usepackage[margin=25mm]{geometry}
\usepackage{fontspec}
\setmainfont{Latin Modern Roman}
\usepackage{amsmath,amssymb,mathtools,unicode-math}
\setmathfont{Latin Modern Math}
\usepackage{enumitem,needspace,xurl,hyperref}
\setlength{\parindent}{0pt}
\setlength{\parskip}{5pt}
\setlength{\emergencystretch}{3em}
\newcommand{\N}{\mathbb{N}}
\newcommand{\Z}{\mathbb{Z}}
\newcommand{\Q}{\mathbb{Q}}
\newcommand{\R}{\mathbb{R}}
\newcommand{\C}{\mathbb{C}}
\newcommand{\F}{\mathbb{F}}
\newcommand{\A}{\mathbb{A}}
\newcommand{\PP}{\mathbb{P}}
\newcommand{\E}{\mathbb{E}}
\newcommand{\eps}{\varepsilon}
\newcommand{\dd}{\,\mathrm{d}}
\begin{document}
\section*{UnsolvedMath ID: 30000015}
\textbf{Source problem code:} catalogue entry
\subsection{Definitions and mathematical statement}
The source asks for stronger moment-method bounds for large random band matrices, sufficient to determine the local scale of eigenvalue spacings at a spectral edge. A notation for the intended problem, without specifying an absent ensemble, is as follows. Write \(H_{N,b}\) for a random self-adjoint matrix of size \(N\), with bandwidth parameter \(b=b(N)\), and order its eigenvalues as \(\lambda_1\leq\cdots\leq\lambda_N\). The band geometry, entry distributions, normalization, and growth regime of \(b\) are not supplied. When the relevant moments exist, write
\[
M_{2p}(N,b)=\E\operatorname{Tr}(H_{N,b}^{2p}),\qquad p\in\N.
\]
The stated research task is to sharpen moment-method inequalities in the intended band ensemble, with enough control as \(N\) and the moment order grow to identify the scale of consecutive eigenvalue gaps near the relevant limiting spectral edge. If the upper edge is denoted by \(E_+\), the quantities of interest include gaps \(\lambda_{N-j+1}-\lambda_{N-j}\) for eigenvalues near \(E_+\); this notation does not prescribe a particular limiting law or bandwidth regime.

No explicit proposed inequality, admissible moment range, asymptotic edge, or target spacing scale is reproduced in the record. Accordingly, this is a precise account of the supplied research direction, rather than a fully quantified conjecture about an invented random-matrix ensemble. The supplied literature assessment, checked 2026-08-21, says that the sharper inequality proposed in the report is not proved in the cited Khorunzhy paper and that no later proof of that exact bound was located in that check. The inequality itself must be recovered from the cited sources before an exact bound can be formalized.
\subsection{Short English statement}
Improve moment estimates for large random band matrices enough to identify eigenvalue spacing near a spectral edge. The record omits the ensemble, bandwidth regime, and proposed sharper inequality, so it does not determine a unique quantified conjecture.
\subsection{Sources}
\begin{itemize}
\item Local UnsolvedMath registry, permanent ID 30000015, catalogue entry, \emph{Spectral-Edge Spacing in Large Random Band Matrices}.
\item Supplied report citation: \emph{Mathematics and Physics of Disordered Systems} (2005), Oberwolfach Reports, pp. 1167--1232, DOI \url{https://doi.org/10.4171/owr/2004/22}. Official report link supplied: \url{https://ems.press/content/serial-article-files/45912?nt=1}.
\item Khorunzhy, supplied preprint link \url{https://arxiv.org/abs/math-ph/0404017}; the record gives no title or theorem number.
\item Supplied literature triage dated 2026-08-21; it lists the Khorunzhy preprint and official report as checked sources.
\end{itemize}
\end{document}
